%
% $Autor: Wings $
% $Datum: 2020-01-29 07:55:27Z $
% $Pfad: Nano33BLESense.tex $
% $Version: 1785 $
%

\chapter{Arduino Nano 33 BLE Sense}

Arduino ist eine Open-Source-Elektronikplattform, die auf leicht zugänglicher Hardware und Software basiert. Sie stellt Mikrocontroller- und Mikroprozessorkits bereit, mit denen sich digitale und analoge Anwendungen realisieren lassen. Der Mikrocontroller ist dabei ein eingebetteter, programmierbarer Rechner, der Sensoren auswerten, Aktoren ansteuern und Daten über Schnittstellen weitergeben kann.

Das Arduino-Team stellt für diese Boards eine einheitliche Entwicklungsumgebung bereit. Dadurch lassen sich Programme in der Arduino IDE oder in kompatiblen Entwicklungsumgebungen schreiben, kompilieren und über USB auf das Board übertragen. Viele hardwarenahe Details werden durch Bibliotheken abstrahiert, sodass Sensoren, Kommunikationsschnittstellen und Peripherie mit vergleichsweise wenig Quellcode genutzt werden können.

\section{Einordnung des Arduino Nano 33 BLE Sense}

Die Arduino-Nano-Familie umfasst kompakte Entwicklungsboards im klassischen Nano-Formfaktor. Der Arduino Nano 33 BLE Sense ist ein besonders leistungsfähiges Modell dieser Serie, da er Bluetooth Low Energy, einen schnellen 32-bit-Mikrocontroller und mehrere bereits integrierte Sensoren kombiniert. Die Bezeichnung „Nano“ beschreibt die kompakte Bauform, „33“ verweist auf die Betriebsspannung von 3,3~V, „BLE“ kennzeichnet die integrierte Bluetooth-Low-Energy-Funktionalität und „Sense“ steht für die umfangreiche On-Board-Sensorik.

Als zentrale Recheneinheit kommt ein Nordic nRF52840 im NINA-B306-Modul zum Einsatz. Dieser Mikrocontroller enthält einen Arm Cortex-M4F-Prozessor mit 64~MHz Taktfrequenz, 1~MB Flash-Speicher und 256~KB SRAM. Durch die integrierte Gleitkommaeinheit eignet sich das Board auch für rechenintensivere Sensoranwendungen, Signalverarbeitung und einfache Machine-Learning-Anwendungen am Rand des Netzwerks.

Das Board arbeitet ausschließlich mit 3,3-V-Logikpegeln. Die Ein- und Ausgänge sind nicht 5-V-tolerant. Deshalb dürfen externe Sensoren oder Aktoren, die mit 5~V arbeiten, nicht direkt mit den I/O-Pins verbunden werden. Bei 5-V-Signalen müssen Pegelwandler, Spannungsteiler oder geeignete 3,3-V-kompatible Module verwendet werden.

\section{Arduino Nano 33 BLE Sense, Sense Lite und Rev. 2}

Für eine technische Dokumentation ist es wichtig, die verschiedenen Varianten des Boards sauber voneinander zu unterscheiden. Der Arduino Nano 33 BLE Sense der ersten Revision besitzt unter anderem die Sensoren LSM9DS1, LPS22HB, HTS221, APDS-9960 und MP34DT05-A. Die Lite-Variante wurde im Zusammenhang mit dem Tiny Machine Learning Kit ausgeliefert und unterscheidet sich im Wesentlichen dadurch, dass der HTS221-Sensor für relative Luftfeuchtigkeit und Temperatur nicht bestückt ist. Dadurch kann die Lite-Version keine relative Luftfeuchtigkeit messen. Eine Temperaturmessung ist weiterhin über den LPS22HB-Drucksensor möglich, allerdings ersetzt dies keine Feuchtemessung.

Der Arduino Nano 33 BLE Sense Rev. 2 verwendet teilweise andere Sensoren als die erste Revision. Bei Rev. 2 wurde die IMU LSM9DS1 durch die Kombination aus BMI270 und BMM150 ersetzt. Außerdem wurde der Temperatur- und Feuchtigkeitssensor HTS221 durch den HS3003 und das Mikrofon MP34DT05 durch das MP34DT06JTR ersetzt. Deshalb müssen bei Rev.-2-Projekten auch die passenden Bibliotheken verwendet werden. Quellcode, der direkt auf die alten Sensorbibliotheken zugreift, ist nicht ohne Anpassung auf Rev. 2 übertragbar.

\section{Arduino Tiny Machine Learning Kit}

Das Arduino Tiny Machine Learning Kit ist eine Entwicklungsplattform für eingebettete Machine-Learning-Anwendungen. Es enthält ein Arduino Nano 33 BLE Sense Lite Board, ein Kameramodul, das Arduino Tiny Machine Learning Shield und ein USB-Verbindungskabel. In diesem Kit dient das Nano-Board als kompakte Recheneinheit, auf der kleine Modelle mit TensorFlow Lite for Microcontrollers oder vergleichbaren Frameworks direkt auf dem Mikrocontroller ausgeführt werden können.

Enthalten sind typischerweise:
\begin{itemize}
	\item Arduino Nano 33 BLE Sense Lite,
	\item Kameramodul,
	\item Arduino Tiny Machine Learning Shield,
	\item USB-Verbindungskabel.
\end{itemize}

\section{Technische Kenndaten}

\begin{table}[h]
	\centering
	\begin{tabular}{ll}
		\textbf{Eigenschaft} & \textbf{Wert} \\
		\hline
		Mikrocontroller & Nordic nRF52840 \\
		Prozessorkern & Arm Cortex-M4F mit FPU \\
		Taktfrequenz & 64~MHz \\
		Betriebsspannung & 3,3~V \\
		I/O-Pegel & 3,3~V, nicht 5-V-tolerant \\
		Flash-Speicher & 1~MB \\
		SRAM & 256~KB \\
		Digitale I/O-Pins & 14 \\
		Analoge Eingänge & 8, 12-bit ADC \\
		PWM & über digitale Pins möglich \\
		Schnittstellen & UART, SPI, I\textsuperscript{2}C, USB, BLE \\
		Abmessungen & ca. 45~mm $\times$ 18~mm \\
	\end{tabular}
	\caption{Technische Grunddaten des Arduino Nano 33 BLE Sense}
	\label{tab:Nano33BLEGrunddaten}
\end{table}

Das Board besitzt kein integriertes Ladesystem für Akkus. Die Versorgung erfolgt über USB oder über die Versorgungsanschlüsse an den Stiftleisten. Der 5-V-Pin darf nicht mit einem klassischen 5-V-Ausgang älterer Arduino-Nano-Boards verwechselt werden. Er ist nur unter bestimmten Bedingungen mit der USB-Versorgung verbunden und eignet sich nicht als allgemeine 5-V-Versorgung für externe Schaltungen.

\section{Pinbelegung und elektrische Hinweise}

Die Pins des Arduino Nano 33 BLE Sense können für digitale Ein- und Ausgänge, analoge Messungen und serielle Schnittstellen genutzt werden. Die analogen Eingänge A0 bis A7 messen Spannungen im Bereich von 0 bis 3,3~V. Die digitalen Pins erkennen ebenfalls 3,3~V als High-Pegel. Eine direkte Verbindung mit 5-V-Ausgängen, beispielsweise von älteren Sensorboards oder klassischen Arduino-Boards, kann das Nano 33 BLE Sense beschädigen.

\begin{longtable}{|c|p{2.5cm}|p{3cm}|p{6.5cm}|}
	\caption{Pinbelegung des Arduino Nano 33 BLE Sense}
	\label{tab:Pinbelegung} \\
	\hline
	\textbf{Pin} & \textbf{Bezeichnung} & \textbf{Typ} & \textbf{Beschreibung} \\
	\hline
	\endfirsthead
	\hline
	\textbf{Pin} & \textbf{Bezeichnung} & \textbf{Typ} & \textbf{Beschreibung} \\
	\hline
	\endhead
	1 & D13 & Digital & GPIO, integrierte orange LED \\
	\hline
	2 & +3V3 & Power Out & Geregelte 3,3-V-Ausgabe für externe 3,3-V-Komponenten \\
	\hline
	3 & AREF & Analog & Analoge Referenzspannung, abhängig von der verwendeten Konfiguration \\
	\hline
	4--11 & A0--A7 & Analog/Digital & Analoge Eingänge, auch als GPIO nutzbar; A4/A5 zusätzlich I\textsuperscript{2}C \\
	\hline
	12 & VUSB & Power & USB-Versorgung, abhängig von der Lötbrücke beziehungsweise Board-Revision \\
	\hline
	13, 18 & RST & Digital In & Aktiv-Low-Reset-Eingang \\
	\hline
	14, 19 & GND & Power & Masseanschluss \\
	\hline
	15 & VIN & Power In & Eingang für externe Versorgung im zulässigen Bereich \\
	\hline
	16--17 & TX/RX & Digital & UART-Schnittstelle, auch als GPIO nutzbar \\
	\hline
	20--28 & D2--D10 & Digital & GPIO, mehrere Pins unterstützen PWM \\
	\hline
	29--30 & D11/D12 & Digital/SPI & SPI MOSI und MISO, auch als GPIO nutzbar \\
	\hline
\end{longtable}

\section{Verbaute Sensoren}

\subsection{APDS-9960: Gesten-, Näherungs-, Farb- und Umgebungslichtsensor}

Der APDS-9960 ist ein digitaler Sensor zur Erfassung von Gesten, Näherung, RGB-Farbe und Umgebungslicht. Er nutzt optische Messverfahren und integrierte Photodioden, um reflektiertes Licht auszuwerten. Dadurch kann das Board beispielsweise Handbewegungen vor dem Sensor erkennen, Abstände im Nahbereich abschätzen oder die Farbe beziehungsweise Helligkeit der Umgebung erfassen.

Typische Anwendungen sind einfache Gestensteuerungen, automatische Helligkeitsanpassung, Farberkennung und berührungslose Bedienkonzepte. Für Arduino-Projekte wird der Sensor üblicherweise über die Bibliothek \texttt{Arduino\_APDS9960} angesprochen.

\subsection{LSM9DS1: Inertialsensor der ersten Revision}

Der LSM9DS1 ist eine 9-Achsen-IMU und kombiniert Beschleunigungssensor, Gyroskop und Magnetometer in einem Baustein. Der Beschleunigungssensor misst lineare Beschleunigungen entlang dreier Achsen. Das Gyroskop misst Winkelgeschwindigkeiten, während das Magnetometer das Magnetfeld erfasst und zum Beispiel für eine Kompassfunktion genutzt werden kann.

IMU-Daten werden häufig für Bewegungserkennung, Lagebestimmung, einfache Navigation, Gestenklassifikation oder TinyML-Anwendungen verwendet. Dabei muss beachtet werden, dass Rohdaten verrauscht sein können und insbesondere Gyroskopdaten über die Zeit driften. Drift bedeutet, dass kleine Messfehler durch Integration zu immer größeren Winkelfehlern anwachsen können. Deshalb sollten IMU-Anwendungen Kalibrierung, Filterung und Plausibilitätsprüfungen enthalten.

\subsection{BMI270 und BMM150: IMU der Rev. 2}

Beim Arduino Nano 33 BLE Sense Rev. 2 wurde der LSM9DS1 durch zwei Sensoren ersetzt. Der BMI270 stellt Beschleunigungs- und Gyroskopdaten bereit, während der BMM150 das Magnetfeld misst. Funktional ergibt sich dadurch ebenfalls eine 9-Achsen-Bewegungserfassung, die Programmierung erfolgt jedoch über andere Bibliotheken als bei der ersten Revision.

\subsection{LPS22HB: Druck- und Temperatursensor}

Der LPS22HB ist ein digitaler barometrischer Drucksensor. Er misst den absoluten Luftdruck und kann zusätzlich Temperaturdaten bereitstellen, die vor allem zur Kompensation der Druckmessung dienen. Aus dem Luftdruck kann näherungsweise eine Höhenänderung berechnet werden, weshalb der Sensor für einfache Wetterstationen, Höhenmessungen und Umgebungsüberwachung geeignet ist.

Der typische Druckmessbereich liegt zwischen 260~hPa und 1260~hPa. Die Daten werden digital über I\textsuperscript{2}C oder SPI übertragen. In Arduino-Projekten wird häufig die Bibliothek \texttt{Arduino\_LPS22HB} verwendet.

\subsection{HTS221 beziehungsweise HS3003: Feuchte- und Temperatursensor}

Der HTS221 ist beim Arduino Nano 33 BLE Sense der ersten Revision für relative Luftfeuchtigkeit und Temperatur zuständig. Beim Arduino Nano 33 BLE Sense Rev. 2 wird diese Aufgabe vom HS3003 übernommen. Die Sense-Lite-Variante besitzt den HTS221 nicht, weshalb sie keine relative Luftfeuchtigkeit messen kann.

Feuchte- und Temperatursensoren eignen sich für Raumklimaüberwachung, Wetterstationen, Smart-Home-Anwendungen und Umweltmessungen. Bei Messungen muss berücksichtigt werden, dass die Position des Boards und Eigenerwärmung die Temperaturwerte beeinflussen können. Für präzise Messungen sollte das Board ausreichend belüftet montiert und nicht direkt neben wärmeerzeugenden Komponenten betrieben werden.

\subsection{MP34DT05-A beziehungsweise MP34DT06JTR: Digitales Mikrofon}

Das integrierte Mikrofon nimmt Schall als digitales PDM-Signal auf. Beim Arduino Nano 33 BLE Sense der ersten Revision wird das MP34DT05-A verwendet, während Rev. 2 ein MP34DT06JTR einsetzt. Das Mikrofon eignet sich für Lautstärkeerfassung, einfache Spracherkennung, Geräuscherkennung und TinyML-Klassifikationen.

Der Zugriff erfolgt über die PDM-Bibliothek. PDM steht für Pulse Density Modulation. Dabei wird ein analoges Audiosignal in einen hochfrequenten digitalen Bitstrom umgewandelt, aus dem durch Verarbeitung wieder Audiodaten gewonnen werden können.

\subsection{ATECC608A: Kryptografischer Co-Prozessor}

Der ATECC608A ist ein kryptografischer Co-Prozessor zur sicheren Speicherung von Schlüsseln, Zertifikaten und kryptografischen Daten. Er kann für Authentifizierung, sichere Kommunikation und IoT-Anwendungen eingesetzt werden. Für einfache Sensorprojekte ist dieser Baustein nicht zwingend erforderlich, wird aber bei vernetzten Anwendungen relevant.

\section{Bibliotheken und Installation}

Die Sensoren des Arduino Nano 33 BLE Sense werden über spezifische Arduino-Bibliotheken angesprochen. Die Installation erfolgt in der Arduino IDE über \texttt{Sketch > Include Library > Manage Libraries}. Dort können die benötigten Bibliotheken gesucht und installiert werden.

Für die erste Revision sind besonders folgende Bibliotheken relevant:
\begin{itemize}
	\item \texttt{Arduino\_LSM9DS1} für Beschleunigung, Gyroskop und Magnetometer,
	\item \texttt{Arduino\_LPS22HB} für Luftdruck und Temperatur,
	\item \texttt{Arduino\_HTS221} für Temperatur und relative Luftfeuchtigkeit,
	\item \texttt{Arduino\_APDS9960} für Gesten, Näherung, Farbe und Umgebungslicht,
	\item \texttt{PDM} für das digitale Mikrofon,
	\item \texttt{ArduinoBLE} für Bluetooth Low Energy.
\end{itemize}

Für Rev. 2 müssen teilweise andere Bibliotheken verwendet werden:
\begin{itemize}
	\item \texttt{Arduino\_BMI270\_BMM150} für die IMU,
	\item \texttt{Arduino\_HS300x} für Temperatur und relative Luftfeuchtigkeit,
	\item \texttt{PDM} für das digitale Mikrofon,
	\item \texttt{ArduinoBLE} für Bluetooth Low Energy.
\end{itemize}

Bei der Sense-Lite-Variante darf keine Feuchtemessung mit dem HTS221 vorausgesetzt werden. Wenn relative Luftfeuchtigkeit benötigt wird, muss ein externer Feuchtigkeitssensor angeschlossen werden.

\section{Beispiel: Zugriff auf das digitale Mikrofon}

Das folgende Beispiel liest Audiodaten über die PDM-Bibliothek ein und gibt einen einfachen Lautstärkewert über die serielle Schnittstelle aus.

\begin{code}
\begin{Arduino}
// Include the PDM library
#include <PDM.h>

// Buffer to store the microphone data
short sampleBuffer[256];

// Variable to store the sound level
int soundLevel = 0;

// Callback function for PDM data
void onPDMdata() {
	int bytesAvailable = PDM.available();
	PDM.read(sampleBuffer, bytesAvailable);

	soundLevel = 0;
	for (int i = 0; i < bytesAvailable / 2; i++) {
		soundLevel += abs(sampleBuffer[i]);
	}
	soundLevel /= bytesAvailable / 2;
}

void setup() {
	Serial.begin(9600);
	while (!Serial);

	PDM.onReceive(onPDMdata);
	PDM.begin(1, 16000);
}

void loop() {
	Serial.println(soundLevel);
	delay(100);
}
\end{Arduino}
\caption{Einfaches Beispiel zur Auswertung des integrierten Mikrofons}
\label{code:microphone}
\end{code}

Der Wert \texttt{soundLevel} ist kein kalibrierter Schalldruckpegel in dB, sondern ein einfacher relativer Messwert. Für präzise akustische Messungen wären Kalibrierung, Fensterung, Filterung und eine definierte Umrechnung erforderlich.

\section{Kalibrierung, Drift und Messgenauigkeit}

Sensoren liefern nicht automatisch ideale Messwerte. Messwerte können durch Rauschen, Offsetfehler, Temperatureinflüsse, Einbaulage und elektrische Störungen beeinflusst werden. Deshalb sollten Sensorwerte vor der weiteren Verarbeitung geprüft, gefiltert und bei Bedarf kalibriert werden.

Bei einer IMU ist die Kalibrierung besonders wichtig. Beschleunigungssensoren können Offset- und Skalierungsfehler besitzen. Gyroskope können einen Nullpunktfehler aufweisen, der bei der Winkelberechnung zu Drift führt. Magnetometer werden zusätzlich durch harte und weiche magnetische Störungen beeinflusst. Harte Störungen entstehen beispielsweise durch dauerhaft magnetisierte Bauteile, weiche Störungen durch Materialien, die das Magnetfeld lokal verzerren.

Eine einfache Kalibrierung besteht darin, den Sensor in verschiedene definierte Lagen zu bringen und Minima, Maxima und Mittelwerte der Messachsen zu bestimmen. Für Magnetometer können Verfahren wie die Ellipsoid-Kalibrierung genutzt werden, bei denen Offset, Skalierung und Achsverkopplungen korrigiert werden. Nach der Kalibrierung sollten die Rohdaten und die korrigierten Werte dokumentiert werden, damit die Messqualität nachvollziehbar bleibt.

\section{Laufzeit, Speicherbedarf und Leistungsaufnahme}

Bei Mikrocontrollerprojekten müssen Laufzeit, Speicherbedarf und Energieverbrauch berücksichtigt werden. Der nRF52840 bietet 1~MB Flash und 256~KB SRAM. Für einfache Sensorprogramme ist das ausreichend. Bei TinyML-Anwendungen können Modellgröße, Tensor-Arena, Puffer und Bibliotheken den verfügbaren Speicher jedoch schnell begrenzen.

Die Laufzeit eines batteriebetriebenen Projekts hängt von der Stromaufnahme des Boards, der Abtastrate der Sensoren, der Bluetooth-Nutzung und dem verwendeten Energiesparmodus ab. Dauerhaft aktive Sensoren, serielle Ausgabe und BLE-Kommunikation erhöhen den Verbrauch. Für lange Laufzeiten sollten Messintervalle reduziert, nicht benötigte Sensoren deaktiviert und Schlafmodi genutzt werden.

\section{Bluetooth Low Energy mit dem NINA-B306-Modul}

Die drahtlose Kommunikation erfolgt über das NINA-B306-Modul mit Bluetooth Low Energy. BLE eignet sich für energieeffiziente Datenübertragung über kurze Distanzen, beispielsweise zur Verbindung mit einem Smartphone oder einem Gateway. In Arduino-Projekten wird dafür typischerweise die Bibliothek \texttt{ArduinoBLE} verwendet.

Ein mögliches Anwendungsszenario ist das Bereitstellen von Sensordaten als BLE-Service. Dabei definiert das Programm Services und Characteristics, über die ein Smartphone Messwerte lesen oder Konfigurationswerte schreiben kann. Für Tests kann beispielsweise die App \textit{nRF Connect} verwendet werden.

\section{Software-Dokumentation mit Doxygen}

Für eine technische Projektarbeit sollte der Quellcode nicht nur funktionieren, sondern auch nachvollziehbar dokumentiert sein. Doxygen eignet sich dafür, aus strukturierten Kommentaren automatisch eine Code-Dokumentation zu erzeugen. Besonders wichtig sind Kommentare für Dateien, globale Konstanten, Funktionen, Klassen und wichtige Datenstrukturen.

Ein typischer Funktionskommentar beschreibt Zweck, Parameter, Rückgabewert und besondere Randbedingungen:

\begin{code}
\begin{Arduino}
/**
 * @brief Reads the current pressure value from the LPS22HB sensor.
 *
 * @return Pressure in hPa. Returns -1.0 if the sensor is not available.
 */
float readPressure()
{
	if (!BARO.begin()) {
		return -1.0;
	}
	return BARO.readPressure();
}
\end{Arduino}
\caption{Beispiel für einen Doxygen-kompatiblen Funktionskommentar}
\label{code:doxygen}
\end{code}

Neben der Quellcodedokumentation sollten auch der Programmablauf, verwendete Bibliotheken, Pinbelegung, Messablauf und bekannte Einschränkungen beschrieben werden. Sinnvoll sind außerdem Ablaufdiagramme, Blockdiagramme und eine kurze Anleitung zur Inbetriebnahme.

\section{Interpretation der Messergebnisse}

Messwerte sollten nicht nur ausgegeben, sondern auch interpretiert werden. Dazu gehört, die physikalische Bedeutung des Messwerts, die Einheit, den erwarteten Wertebereich und mögliche Fehlerquellen zu erklären. Bei Temperatur- und Feuchtemessungen können beispielsweise Eigenerwärmung und Luftzirkulation eine Rolle spielen. Bei IMU-Daten müssen Lage, Kalibrierung, Drift und Filterung berücksichtigt werden. Bei Audiodaten ist zu unterscheiden, ob nur ein relativer Lautstärkewert oder ein kalibrierter Schalldruckpegel angegeben wird.

Für die Dokumentation empfiehlt es sich, Messwerte in Tabellen oder Diagrammen darzustellen und kurz zu bewerten. Dabei sollte angegeben werden, unter welchen Bedingungen gemessen wurde, welche Bibliothek verwendet wurde und ob die Werte gefiltert oder kalibriert wurden.

\section{Literatur und Abbildungen}

Für die technische Dokumentation sollten alle Datenblätter, offiziellen Arduino-Seiten, Bibliotheksdokumentationen und verwendeten Abbildungen vollständig angegeben werden. Abbildungen benötigen eine aussagekräftige Bildunterschrift, eine eindeutige Quellenangabe und einen Bezug im Fließtext. Eigene Schaltpläne, Pinbelegungen und Blockdiagramme sollten als eigene Abbildungen gekennzeichnet werden.

Wenn Bilder aus offiziellen Datenblättern oder Webseiten übernommen werden, muss geprüft werden, ob die Nutzung im Rahmen der Arbeit erlaubt ist. Alternativ können eigene TikZ-Abbildungen, Fritzing-Schaltpläne oder selbst erstellte Blockdiagramme verwendet werden.
